\documentclass[10pt,a4paper]{amsart}
\usepackage[margin=19mm]{geometry}
\usepackage{amsmath,amssymb,mathtools}
\usepackage[scr=rsfs,cal=euler]{mathalfa}
\usepackage{xfrac}
\usepackage[unicode,hidelinks,bookmarks=false]{hyperref}
\newcommand{\ie}{i.e.\ }
\newcommand{\cf}{cf.\ }
\newcommand{\resp}{respectively\ }
\newcommand{\Subsection}{\subsection}
\providecommand{\nobreakdash}{\nobreak\hbox{-}}
\setlength{\emergencystretch}{2em}
\multlinegap0pt
\usepackage{bm}
\usepackage{bbm}
\usepackage{dsfont}
\usepackage{xspace}
\usepackage{cancel}
\usepackage{mathtools}
\usepackage{xcolor}
\usepackage{amsthm}
\makeatletter
\@ifundefined{theorem}{\newtheorem{theorem}{Theorem}}{}
\@ifundefined{lemma}{\newtheorem{lemma}[theorem]{Lemma}}{}
\@ifundefined{proposition}{\newtheorem{proposition}[theorem]{Proposition}}{}
\makeatother
\newcommand{\R}{\mathbb{R}}
\newcommand{\Z}{\mathbb{Z}}
\title[The Lyapunov spectrum for Schneider map on $p\mathbb{Z}\_p$]{The Lyapunov spectrum for Schneider map on $p\mathbb{Z}\_p$}
\author{Matias Alvarado, Nicol\'as Ar\'evalo-Hurtado}
\date{Source: arXiv:2601.05915v1 (CC BY 4.0)}
\begin{document}
\maketitle

\noindent\emph{Source and licence.} The Lyapunov spectrum for Schneider map on $p\mathbb{Z}_p$. Version arXiv:2601.05915v1 (CC BY 4.0), distributed under the Creative Commons Attribution 4.0 International licence, as recorded in the arXiv licence metadata for this exact version and shown on the abstract page https://arxiv.org/abs/2601.05915v1. Full rights notice: licenses/arxiv-2601.05915-CC-BY-4.0.txt.\par\smallskip

\noindent\textit{Editorial scope.} Counted unit: the Proposition whose source label is \texttt{Proposition 1} in the article (source0.tex lines 1136--1141, source label \texttt{Proposition 1}), delivered together with the article's own complete original proof of it (source0.tex lines 1142--1204, one \texttt{proof} environment of 63 lines). No mathematics is written, repaired or re-derived: statement and proof are the article's own text line for line. Required context retained uncounted: Lemma3 (lines 1077--1084). Every definition, display and label of those lines is retained unchanged. Omitted from this package and not redistributed: 1--1076, 1085--1135, 1205--1265. Nothing inside a retained excerpt is omitted or altered: every retained body is delivered exactly as the article's lines are written, so no omission receipt of any kind accompanies this package. Citations: the retained excerpts carry no citation command at all, so no bibliography entry of the article is transcribed and no reference is redistributed. Reference adaptation: none was needed and none was made; every reference inside the delivered unit resolves inside it, so no unresolved reference or citation is produced. Printed slips kept literal: no printed word, symbol, hypothesis or number of the counted statement or of its proof was altered. Layout and class: the article's own class is \texttt{amsart} with packages including amsmath, mathtools, fullpage, xspace, amssymb, inputenc, graphicx, color, hyperref, xcolor, amsthm, amscd, amsfonts, stmaryrd; the wrapper is the lane's pinned \texttt{amsart} editorial wrapper at 10pt on A4, which restates the theorem-like environments the retained text needs and adds the article's own macro definitions that the retained text needs (\texttt{\textbackslash R}, \texttt{\textbackslash Z}), with extra wrapper packages (bm, bbm, dsfont, xspace, cancel, mathtools, xcolor). The delivered numbering is the unit's own. Assets: no retained span contains a figure environment, a graphics-inclusion command, a PDF-page inclusion, a table or a TikZ picture; no image file is copied into this package, the package declares no asset dependency and no image file is redistributed with it. Licence: version arXiv:2601.05915v1 of this article is distributed under the Creative Commons Attribution 4.0 International licence, as recorded in the arXiv licence metadata for this exact version and shown on the abstract page https://arxiv.org/abs/2601.05915v1. A full rights notice is delivered at licenses/arxiv-2601.05915-CC-BY-4.0.txt. Characters: every retained excerpt is pure ASCII, so the delivered engine, XeLaTeX with its default Latin Modern text font, reports no missing character.
\begin{lemma}\label{Lemma3}
    For each $q\in \R$ we obtain
    \begin{enumerate}
        \item $\nu_{q}(\pi_{n}(\tilde{K}_{\alpha(q)}))=1$.
        \item $d_{\nu_{q}}(x)=\mathcal{T}(q)+q\alpha(q)$ for $\nu_{q}$-almost every point $x\in \pi_{n}(\tilde{K}_{\alpha(q)})$.
        \item $\dim_{\mathrm{H}}\pi_{n}(\tilde{K}_{\alpha(q)})=\mathcal{T}(q)+q\alpha(q)$.
    \end{enumerate}
\end{lemma}

\begin{proposition}\label{Proposition 1}
\hfill    \begin{enumerate}
        \item For each $q\in \mathbb{R}$, $w\in \tilde{K}_{\alpha(q)}$, and $x\in \pi_{n}(w)$, we have $d_{\nu_{q}}(x)=\alpha(q)$.
        \item For each $q\in \mathbb{R}$ and each $x\in K_{\alpha(q)}$ there exists $w\in \tilde{K}_{\alpha(q)}$ such that $\pi_{n}(w)=x$.  In other words, for each $q\in \mathbb{R}$ we have $\pi_{n}(\tilde{K}_{\alpha(q)})=K_{\alpha(q)}$. 
    \end{enumerate}
\end{proposition}

\begin{proof}
     Recall first that for $x,y\in  p\Z_{p,n}\setminus F$ such that $|x-y|_{p}<r_{0}$ with  $r_{0}=p^{-(n+1)}$ we have 
     \begin{align*}
         |T_{p,n}(x)-T_{p,n}(y)|_{p}\geq p|x-y|_{p}.
     \end{align*}Indeed, for any $0<r<r_{0}$, there exist $N(r)>0$ and positive constants $C_{12}$ and $C_{13}$ such that if $0\leq m \leq N(r)$, then for all $x\in  p\Z_{p,n}\setminus F$
     \begin{align}\label{Ecu10}
         C_{12}r\prod^{m-1}_{k=0}\psi(T_{p,n}^{k}(x))\leq  |T_{p,n}^{m}(B(x,r))|_{p}\leq C_{13}r\prod^{m-1}_{k=0}\psi(T^{k}_{p,n}(x)).
     \end{align} This follows from
     \begin{align*}
         |T_{p,n}^{m}(B(x,r))|_{p}\geq p^{a_{1}+\cdots +a_{m}}r\geq p^{m}r\dfrac{\prod^{m-1}_{k=0}\psi(T_{p,n}^{k}(x))}{\prod^{m-1}_{k=0}\psi(T_{p,n}^{k}(x))}\geq C_{12}r\prod^{m-1}_{k=0}\psi(T_{p,n}^{k}(x)),
     \end{align*}where $C_{12}=p^{m(1-n)}$, and
     \begin{align*}
         |T_{p,n}^{m}(B(x,r))|_{p}\leq 1\cdot \dfrac{r\prod^{m-1}_{k=0}\psi(T_{p,n}^{k}(x))}{r\prod^{m-1}_{k=0}\psi(T_{p,n}^{k}(x))}\leq C_{13}r\prod^{m-1}_{k=0}\psi(T_{p,n}^{k}(x)),
     \end{align*} where $C_{13}=r^{-1}p^{-m}$.  Analogously, there are positive constants $C_{14}$ and $C_{15}$ such that for all $x\in p\Z_{p,n}\setminus F$ and all $(a,b)\in E_{n}$ 
\begin{align}\label{Ecu11}
         C_{14}r\prod^{m-1}_{k=0}\psi(T_{p,n}^{k}(x))^{-1}\leq  |h_{a,b}^{m}(B(x,r))|_{p}\leq C_{15}r\prod^{m-1}_{k=0}\psi(T^{k}_{p,n}(x))^{-1},
     \end{align} where $h_{a,b}$  is the inverse branch of  $T^{-m}_{p,n}$ with image in $I_{a,b}$. Let fix $x\in p\Z p_{n}\setminus F$ and define $N=N(x,r)$ such that 
     \begin{align*}
         C_{13}r\prod^{N-1}_{k=0}\psi(T^{k}_{p,n}(x))\leq r_{0}\leq C_{13}r\prod^{N}_{k=0}\psi(T^{k}_{p,n}(x)).
     \end{align*} So $N(x,r)$ is the number of necessary iterations of $B(x,r)$ in order to growth and still having diameter lower than $r_{0}$. Thus, by equations \eqref{Ecu10} and \eqref{Ecu11} we obtain that
\begin{align*}
    \left|T_{p,n}^{N}(B(x,r))\right|_{p}\leq C_{13}r\prod^{N-1}_{k=0}\psi(T^{k}_{p,n}(x))\leq r_{0}=\left|B(T^{N}_{p,n}(x),r_{0})\right|_{p}.
\end{align*}So $B(x,r)\subset h(B(T^{N}_{p,n}(x),r_{0}))$ where $h$ is the inverse branch of $T^{-N}_{p,n}$ corresponding to $x$. Also, from equation \eqref{Ecu11}
\begin{align*}
    \left|h(B(T^{N}_{p,n}(x),r_{0}))\right|_{p}\leq C_{15}r_{0}\prod^{N-1}_{k=0}\psi(T^{k}_{p,n}(x))^{-1}\leq C_{16}r=|B(x,C_{16}r)|_{p},
\end{align*}where $C_{16}>0$ is a constant. Thus,
\begin{align*}
    B(x,r)\subset h(B(S^{N}(x),r_{0})) \subset B(x,C_{16}r).
\end{align*} Recall there exists $M_{r_{0}}\in \mathbb{N}$ such that $\mathcal{R}(T^{N}_{p,n}(x))=h_{0}(I_{a,b}),$ where $h_{0}$ is the corresponding branch of $T_{p}^{-M_{r_{0}}}$ \textcolor{red}{ } at $I_{(a,b)}$.  There are positive constants $C_{17}$ and $C_{18}$ such that 
\begin{align*}
    B(T^{N}_{p,n}(x),C_{17}r_{0})\subset \mathcal{R}(T^{N}_{p,n}(x))\subset B(T^{N}_{p,n}(x),C_{18}r_{0}).
\end{align*} This implies 
\begin{align*}
    h(B(T^{N}_{p,n}(x),C_{17}r_{0}))\subset h(\mathcal{R}(T^{N}_{p,n}(x)))\subset h(B(T^{N}_{p,n}(x),C_{18}r_{0})), 
\end{align*}and so 
\begin{align}\label{Ecu12}
    \nu_{q}(B(T^{N}_{p,n}(x),C_{17}r_{0}))\subset \nu_{q}(\mathcal{R}(T^{N}_{p,n}(x)))\subset \nu_{q}(B(T^{N}_{p,n}(x),C_{18}r_{0})).
\end{align} On the other hand, from Lemma \ref{Lemma3} and the fact that $\nu_{q}$ is a Gibbs measure, we obtain that is diametrically regular. Indeed, for $C_{19}=C^{2}p^{-1}e^{2P(\rho)}$ we get for any $r_{1}>0$ and $y\in p\Z_{p,n}\setminus F$ that
\begin{align*}
    \nu_{q}(B(r,pr_{1}))\leq C_{19}\nu_{q}(B(x,r_{1})).
\end{align*}Thus
\begin{align*}
    \nu_{q}(B(T^{N}_{p,n}(x),C_{18}r_{0}))&\leq C_{20}\nu_{q}(B(T^{N}_{p,n}(x),r_{0}))=C_{17}\nu_{q}(h(B(T^{N}_{p,n}(x),r_{0})))\\
    & \leq C_{20}\nu_{q}(B(x,C_{16}r))\\
    &\leq C_{21}\nu_{q}(B(x,r)),
\end{align*}where $C_{20}>0$ and $C_{21}>0$ are constants. Similarly
\begin{align*}
    \nu_{q}(B(T^{N}_{p,n}(x),C_{17}r_{0}))\geq C_{22}\nu_{q}(B(T^{N}_{p,n}(x),r_{0}))=C_{22}\nu_{q}(h(B(T^{N}_{p,n}(x),r_{0})))\geq C_{23}\nu_{q}(B(x,r)),
\end{align*}where $C_{22}>0$ and $C_{23}>0$. Thus, by equation \eqref{Ecu12} 
\begin{align*}
    C_{23}\nu_{q}(B(x,r))\leq \nu_{q}(h(R(T^{N}_{p,n}(x)))) \leq  C_{21}\nu_{q}(B(x,r)).
\end{align*}Since $\nu_{q}$ is a Gibbs measure for $\log \vartheta$ and $h(R(T^{N}_{p,n}(x)))$ is an element of the Moran cover, we get
\begin{align*}
    d_{v}(x)=\lim_{r\to 0}\dfrac{\nu_{q}(B(x,r))}{\log r}=\lim_{N(r)\to \infty}\dfrac{\log\prod^{N(r)-1}_{k=0}\vartheta(\sigma^{k}(w))}{\log \prod^{N(r)-1}_{k=0}\psi (\pi_{n}\circ \sigma^{k}(w))^{-1}},
\end{align*}where $\pi_{n}(w)=x$. Now assume that $d_{v}(x)=\alpha(q)$, we need to show the existence of a subsequential limit $N(r)\to \infty$ that implies the limit for $n\to \infty$.  Consider $r_{k}=p^{-k}$, it follows from the definition of $N(r)$ and the uniformly expansiveness of $T^{N}_{p,n}$
\begin{align*}
    N(r_{k+1})-N(r_{k})=\begin{cases}
        0 & \text{if $N(x,r_{k+1})=N(x,r_{k})$}\\
        1 & \text{if $N(x,r_{k+1})>N(x,r_{k})$}
    \end{cases}.
\end{align*}
And so the proposition follows. %Seguramente las constantes esten en terminos de p^{-n}.
\end{proof}

\end{document}
